\chapter{Introdução e como usar este guia}

\section{Objetivo}
Este guia acompanha o pacote \code{aero\_pursuit} e serve de base teórica e prática
para o TCC: \emph{comparar o desempenho de diferentes leis de controle} na perseguição
de um drone alvo por um x500 equipado com câmera de profundidade (OAK-D Lite) e
controlado por uma Pixhawk/PX4 em modo \emph{offboard}.

Para cada técnica implementada no código (\cref{tab:controladores}) o guia mostra:
\begin{enumerate}
  \item a lei de controle exatamente como está no C++ (arquivo e equação);
  \item o modelo usado no projeto e as hipóteses;
  \item como \textbf{calcular os ganhos} à mão (alocação de polos, especificações de
        tempo de resposta e sobressinal) e pelo \textbf{lugar geométrico das raízes} (LGR);
  \item como obter a versão \textbf{ótima} (LQR/DARE, MPC, otimização numérica de um
        índice de desempenho);
  \item como levar em conta o \textbf{tempo discreto} (período do laço de 50 ms, ZOH),
        a \textbf{malha interna da Pixhawk} e os \textbf{atrasos} (câmera, rastreador,
        comunicação);
  \item o script MATLAB correspondente na pasta \code{matlab/}.
\end{enumerate}

\begin{table}[H]
\centering
\caption{Controladores do pacote (\code{controller\_type} em \code{aero\_pursuit.yaml}).}
\label{tab:controladores}
\begin{tabular}{@{}llll@{}}
\toprule
\code{controller\_type} & Técnica & Saída & Capítulo \\
\midrule
\code{pbvs\_pid} & Servovisão por posição + PID & velocidade & \ref{cap:pid} \\
\code{lqr} & Regulador linear quadrático (DARE) & aceleração & \ref{cap:lqr} \\
\code{lqi} & LQR com ação integral & aceleração & \ref{cap:lqi} \\
\code{asmc} & Modos deslizantes adaptativo & aceleração & \ref{cap:smc} \\
\code{mpc} (\code{kinematic}) & MPC linear LTV (duplo integrador) & aceleração & \ref{cap:mpc} \\
\code{mpc} (\code{dynamic}) & MPC de estado completo do x500 & atitude/empuxo & \ref{cap:mpcdin} \\
\code{hermite} & Planejamento quíntico + planeza & vel./acel./jerk & \ref{cap:hermite} \\
\code{hybrid\_pn} & Guiamento híbrido PN/PP & aceleração & \ref{cap:guiamento} \\
\code{external} & Política externa (RL) & velocidade & --- \\
\bottomrule
\end{tabular}
\end{table}

\section{Mapa do sistema}
O laço de perseguição tem cinco blocos em cascata. Entender quem é rápido e quem é
lento é o ponto de partida de qualquer projeto:

\begin{center}
\small
\begin{tabular}{@{}lll@{}}
\toprule
Bloco & Frequência & Atraso típico \\
\midrule
Câmera RGB/profundidade (Gazebo) & \SI{30}{Hz} & \SI{17}{ms} (meio período) \\
Rastreador VitTrack + profundidade & $\approx$\SI{30}{Hz} & 30--\SI{60}{ms} \\
Estimador (KF de velocidade constante) & por medida & atraso de fase do filtro \\
\textbf{Controlador de perseguição} (\code{pursuit\_controller\_node}) & \SI{20}{Hz} ($T_s=\SI{0,05}{s}$) & $T_s/2$ (ZOH) \\
PX4: velocidade $\to$ atitude $\to$ taxa $\to$ motores & 50/250/\SI{800}{Hz} & $\tau_v\approx$ 0,2--\SI{0,7}{s} \\
\bottomrule
\end{tabular}
\end{center}

O controlador de perseguição é a \emph{malha externa}: ele não comanda motores, e sim
um \code{TrajectorySetpoint} (velocidade + aceleração de \emph{feedforward}) que a
Pixhawk segue com a sua própria malha interna. Portanto a ``planta'' vista pelo nosso
controlador é \textbf{o drone já controlado pela PX4} --- e não o drone nu.

\section{Como ler e usar}
\begin{itemize}
  \item Os capítulos \ref{cap:modelagem}--\ref{cap:discreto} são a base comum: modelagem,
        ferramentas (LGR, Bode, especificações) e o efeito de amostragem e da PX4.
        Leia-os primeiro.
  \item Os capítulos \ref{cap:pid}--\ref{cap:guiamento} são independentes: um por técnica.
  \item O capítulo \ref{cap:comparacao} define a metodologia de comparação justa (mesmos
        cenários, mesmo índice de custo, ajuste ótimo de cada controlador, Monte Carlo).
  \item O capítulo \ref{cap:caso} documenta a correção da oscilação ``vai e volta''
        encontrada no MPC (filtro de Kalman relativo e referência da linha de visada).
        É um bom estudo de caso para o texto do TCC.
  \item O Apêndice \ref{ap:matlab} descreve todos os scripts MATLAB e o gerador dos
        modelos Simulink.
\end{itemize}

\begin{dica}
Rode \code{matlab/parametros.m} primeiro: ele define, num único \code{struct}, todos os
parâmetros (x500, PX4, ganhos do \code{yaml}, atrasos). Todos os outros scripts o
chamam. Se você mudar um ganho no \code{aero\_pursuit.yaml}, mude também lá.
\end{dica}
